% Generated draft: requires mathematical and source review.
\documentclass[11pt]{article}
\usepackage[margin=25mm]{geometry}
\usepackage{fontspec}
\setmainfont{Latin Modern Roman}
\usepackage{amsmath,amssymb,mathtools,unicode-math}
\setmathfont{Latin Modern Math}
\usepackage{enumitem,needspace,xurl,hyperref}
\setlength{\parindent}{0pt}
\setlength{\parskip}{5pt}
\setlength{\emergencystretch}{3em}
\newcommand{\N}{\mathbb{N}}
\newcommand{\Z}{\mathbb{Z}}
\newcommand{\Q}{\mathbb{Q}}
\newcommand{\R}{\mathbb{R}}
\newcommand{\C}{\mathbb{C}}
\newcommand{\F}{\mathbb{F}}
\newcommand{\A}{\mathbb{A}}
\newcommand{\PP}{\mathbb{P}}
\newcommand{\E}{\mathbb{E}}
\newcommand{\eps}{\varepsilon}
\newcommand{\dd}{\,\mathrm{d}}
\begin{document}
\section*{UnsolvedMath ID: 701769}
\textbf{Source problem code:} catalogue entry
\subsection{Definitions and mathematical statement}
A group \(H\) is just infinite if it is infinite and every nontrivial normal subgroup of \(H\) has finite index. A group \(G\) is hereditarily just infinite if every subgroup of finite index in \(G\), including \(G\) itself, is just infinite. It is residually finite if, for every \(g\in G\setminus\{1\}\), there exist a finite group \(Q\) and a homomorphism \(\varphi:G\to Q\) such that \(\varphi(g)\ne 1\). It is linear if there exist a field \(K\), an integer \(n\geq 1\), and an injective homomorphism \(G\to\operatorname{GL}_n(K)\).

For the conventional interpretation of a self-similar action, let \(X\) be a finite alphabet with \(|X|\geq 2\). The rooted tree \(X^*\) consists of all finite words over \(X\), with the empty word as root and edges joining \(w\) to \(wx\) for \(x\in X\). A faithful action \(\rho:G\hookrightarrow\operatorname{Aut}(X^*)\) is self-similar if
\[
\forall g\in G\;\forall x\in X\;\exists h\in G\;\forall w\in X^*,\qquad
\rho(g)(xw)=\rho(g)(x)\,\rho(h)(w),
\]
where juxtaposition on the right denotes concatenation of words.

Under this interpretation, the question is whether there exist a residually finite, hereditarily just infinite group \(G\), a finite alphabet \(X\) with \(|X|\geq 2\), and a faithful self-similar action of \(G\) on \(X^*\), such that
\[
\text{for every field }K\text{ and every integer }n\geq 1,
\quad G\text{ admits no injective homomorphism into }\operatorname{GL}_n(K).
\]
The supplied statement does not explicitly specify faithfulness or the branching conditions on the rooted tree. Thus the formulation above makes the conventional faithful, finite-degree regular-tree interpretation explicit; confirmation of that interpretation from the original problem is required. No finite-generation hypothesis is supplied or imposed.
\subsection{Short English statement}
Does there exist a residually finite group whose every finite-index subgroup is just infinite, which admits a self-similar rooted-tree action and is not linear over any field? The precise formulation above uses a faithful action on a regular rooted tree of finite degree; these action conventions require source confirmation.
\subsection{Sources}
\begin{itemize}
\item Local UnsolvedMath registry, permanent ID \(701769\), exact source problem code \texttt{catalogue entry}, title \emph{Kourovka Notebook Problem 20.111}. The supplied record labels the problem open; this status has not been independently verified.
\item \emph{The Kourovka Notebook (Unsolved Problems in Group Theory)}, Problem 20.111, proposed by M. Vannacci in 2022. The supplied attribution identifies the 20th issue (2022) and the 21st issue edition (2026), with identifier \texttt{arXiv:1401.0300v48}.
\item The supplied collection description attributes editing of the Kourovka Notebook, Issues 1--21 (1965--2026), to E. I. Khukhro and V. D. Mazurov. No source text beyond the supplied registry record has been independently consulted.
\end{itemize}
\end{document}
